% GPT-6 Astra Pro; /root/fixed_surface; 22 November 2026 UTC.
% Effort not exposed. Complete new proof for whole independent review.
\section[HD-F1-S33: coherent sheaves on the actual norm image]{HD-F1-S33: every coherent sheaf on the actual norm image meeting the full S14 character has an invariant trace kernel}
\label{sec:f1-norm-fourfold-coherent-support}

Retain the actual closed immersion $i:Y\hookrightarrow P$,
the quotient algebra $E=i_*\mathcal O_Y$, and the geometric
open of S29. In particular, $Y$ is normal and integral.
Put
\[
 D=\mathfrak D_P,\qquad
 W_E=\operatorname{im}\bigl(
     \operatorname{ob}_E:D\longrightarrow
                    \operatorname{Ext}^2_P(E,E)\bigr).
 \tag{F1S33.1}
\]
The obstruction is the S24 first-jet obstruction, with
the calibrated sign of S31. S29 proves
$\dim W_E=16-r_C\ge4$, where $r_C$ is the rank of the
actual cyclic Prym period differential at the given cover.

Let $\mathcal G$ be an actual coherent sheaf on $Y$, and
form its exact closed-immersion pushforward
\[
 F=i_*\mathcal G,\qquad
 \mathcal O_P\longrightarrow i_*\mathcal O_Y
       \longrightarrow\mathcal End_{\mathcal O_P}(F),
 \qquad I_YF=0.
 \tag{F1S33.2}
\]
For positive generic rank
$r=\dim_{\mathbf C(Y)}\mathcal G_{\operatorname{Spec}\mathbf C(Y)}$,
there is a well-defined surjective linear map
\[
\begin{split}
 \mathcal T_{\mathcal G}:\quad
 V_F&=\operatorname{ev}_F
             \bigl(D\oplus H^2(P,\mathcal O_P)\bigr)
                  \longrightarrow W_E,\\
 \operatorname{ev}_F(\kappa+\eta)
        &\longmapsto-r\,\operatorname{ob}_E(\kappa).
\end{split}
 \tag{F1S33.3}
\]
The domain in (F1S33.3) is the actual characteristic
image in $\operatorname{Ext}^2_P(F,F)$.
Its construction below uses actual module deformations
and traces on the normal support.

Every sheaf (F1S33.2) satisfying the entire S14 requirement
\[
 e^B\operatorname{ch}(F)=\lambda t+Q(\beta),
 \qquad B\in H^{1,1}(P,\mathbf Q),\quad
 0\ne\lambda t\in W_K(P),\quad Q\in\mathbf Q[x]
 \tag{F1S33.4}
\]
has $r>0$. For every specified coherent lift of a faithful
finite point subgroup, it has an invariant full
degree-two trace kernel of dimension at least
$16-r_C\ge4$. Hence these actual sheaves fail weak
point-equivariant semiregularity.

\paragraph{1. Trace on a coherent sheaf over the normal support.}
There is a canonical sheaf morphism
\[
 \operatorname{tr}_{\mathcal G}:
          \mathcal End_Y(\mathcal G)\longrightarrow\mathcal O_Y.
 \tag{F1S33.5}
\]
Here is the construction, including possible torsion
and failure of local freeness of $\mathcal G$.
On an affine normal domain $R\subset K=\mathbf C(Y)$,
the torsion submodule is
$M_{\mathrm{tor}}=\ker(M\to M\otimes_RK)$.
It is finite and preserved by every $R$-linear
endomorphism of $M$.
An endomorphism therefore induces an endomorphism of
the $r$-dimensional $K$-vector space
$(M/M_{\mathrm{tor}})\otimes_RK$.
Take its usual trace in $K$.

For every height-one prime $\mathfrak p$ of $R$,
$R_{\mathfrak p}$ is a discrete valuation ring and
$(M/M_{\mathrm{tor}})_{\mathfrak p}$ is finite free.
The same trace consequently lies in $R_{\mathfrak p}$.
Normality gives the actual intersection identity
\[
             R=\bigcap_{\operatorname{ht}\mathfrak p=1}
                        R_{\mathfrak p}\quad\hbox{inside }K.
 \tag{F1S33.6}
\]
The trace is therefore a regular function on this
affine open. The construction commutes with
localization and glues to (F1S33.5).
Linearity, invariance under conjugation by a sheaf
isomorphism, and
$\operatorname{tr}_{\mathcal G}(1_{\mathcal G})=r$
hold on the generic vector space and hence on $Y$.
This proves the needed extension across the
codimension-at-least-two nonfree locus, as well as
compatibility with torsion.
The precise normal-domain and DVR facts are
Stacks Tags 0AVB, 031S and 0AUW; their complete
displayed proofs were freshly read.

\paragraph{2. A vanishing characteristic class gives a twisted flat sheaf.}
Suppose
\[
       \operatorname{ev}_F(\kappa+\eta)=0,\qquad
       \kappa\in D,\quad \eta\in H^2(P,\mathcal O_P).
 \tag{F1S33.7}
\]
The sheaf $F$ is perfect on the smooth projective
variety $P$. Choose affine opens $U_i$ and
\v Cech representatives $\kappa_{ij}$ and
$\eta_{ijk}$. Let $A_\varepsilon=\mathbf C[\varepsilon]/
(\varepsilon^2)$.
The commutative first-order deformation $P_\kappa$
has local algebras
$\mathcal O_{U_i}[\varepsilon]/(\varepsilon^2)$
with transition maps
\[
 \rho_{ij}(f+\varepsilon h)
       =f+\varepsilon\bigl(h+\kappa_{ij}(f)\bigr).
 \tag{F1S33.8}
\]
The chosen cocycle $\eta$ specifies scalar twisting
with triple product $1-\varepsilon\eta_{ijk}$.

We use the explicit construction in Toda,
\emph{Deformations and Fourier--Mukai transforms},
J.~Differential Geom.~81 (2009), 197--224,
Proposition~6.1 with its Lemmas~6.2--6.4.
Its degree-one convention can be compared directly
on the same local resolution:
\[
 a_{\mathrm{Toda}}(F)=
       \nabla d_F+T(d_{\check C})\nabla
       =\nabla d_F-d_{\check C}\nabla
       =-\operatorname{At}_{\mathrm{jet}}(F).
 \tag{F1S33.9}
\]
Here $T$ is the translation used in that proof,
$T(d_{\check C})=-d_{\check C}$.
The fixed BF characteristic convention of S31 is
$\operatorname{ev}_F(\kappa)
 =-\kappa\mathbin{\lrcorner}\operatorname{At}_{\mathrm{jet}}(F)$.
Consequently the hypothesis in Toda's Proposition~6.1
for $u=(\eta,\kappa,0)$ is exactly
\[
 u\cdot\exp a_{\mathrm{Toda}}(F)
     =\eta\,1_F+\kappa\mathbin{\lrcorner}a_{\mathrm{Toda}}(F)
     =\operatorname{ev}_F(\kappa+\eta)=0.
 \tag{F1S33.10}
\]
Only this directly checked degree-zero and degree-one
comparison is needed.

The cited construction gives an actual perfect
$\eta$-twisted complex $F^\dagger$ on $P_\kappa$
whose derived reduction is $F$.
It also gives its local finite free representatives.
We now extract an actual flat twisted sheaf,
rather than assume that a derived lift is already
a sheaf.
On a sufficiently small open take such a bounded
finite free complex $L^\bullet$ and put
$Q^\bullet=L^\bullet/\varepsilon L^\bullet$.
Termwise flatness gives
\[
 0\longrightarrow\varepsilon L^\bullet
    \longrightarrow L^\bullet\longrightarrow Q^\bullet
    \longrightarrow0,\qquad
 \varepsilon L^\bullet\simeq Q^\bullet.
 \tag{F1S33.11}
\]
The latter isomorphism is induced by multiplication
by $\varepsilon$. The complex $Q^\bullet$ has only
degree-zero cohomology, equal to $F$.
The long exact sequence therefore gives
$H^j(L^\bullet)=0$ for $j\ne0$ and
\[
 0\longrightarrow F
   \longrightarrow M_\varepsilon:=H^0(L^\bullet)
   \longrightarrow F\longrightarrow0.
 \tag{F1S33.12}
\]
Multiplication by $\varepsilon$ on $M_\varepsilon$
is the composite of the quotient and the inclusion
in (F1S33.12). Thus
$\ker(\varepsilon:M_\varepsilon\to M_\varepsilon)
 =\varepsilon M_\varepsilon$, and
$M_\varepsilon/\varepsilon M_\varepsilon=F$.
This is the flatness criterion over $A_\varepsilon$;
it also follows immediately by lifting a vector-space
basis of the central fibre. Hence $M_\varepsilon$
is an actual flat sheaf over $A_\varepsilon$.
Cohomology of the given twisted transition
isomorphisms supplies semilinear isomorphisms
$\phi_{ij}:M_{\varepsilon,i}\to M_{\varepsilon,j}$
with
\[
 \phi_{ki}\phi_{jk}\phi_{ij}
            =(1-\varepsilon\eta_{ijk})\,1.
 \tag{F1S33.13}
\]
The central reductions are the transition
isomorphisms of the given $F=i_*\mathcal G$.

\paragraph{3. The actual ideal-action map.}
Work on one of these opens, with central algebra
$R$, ideal $I=I_Y$, and $B=R/I$.
For $f\in I$, the actual module structure gives
\[
           fM_\varepsilon\subseteq
                      \varepsilon M_\varepsilon\simeq\mathcal G.
 \tag{F1S33.14}
\]
For a local section $m$ of $\mathcal G$, choose a
lift $\widetilde m\in M_\varepsilon$ and define
\[
       a_f(m)=\varepsilon^{-1}(f\widetilde m).
 \tag{F1S33.15}
\]
The inverse symbol means the inverse of the
specified isomorphism
$\mathcal G\xrightarrow{\sim}\varepsilon M_\varepsilon$
in (F1S33.12), not division in the ambient ring.
Replacing $\widetilde m$ by
$\widetilde m+\varepsilon n$ changes
$f\widetilde m$ by $\varepsilon fn=0$, because
$f$ annihilates the central fibre.
Thus (F1S33.15) is independent of the lift.
Commutativity of the ambient algebra gives
$a_f\in\mathcal End_B(\mathcal G)$ and
\[
 a_{hf}=\bar h\,a_f\quad(h\in R),\qquad
 a_{fg}=0\quad(f,g\in I).
 \tag{F1S33.16}
\]
For the second equality, $gM_\varepsilon$ is
already contained in $\varepsilon M_\varepsilon$,
which $f$ annihilates.
Consequently this is an actual $B$-linear
morphism from $I/I^2$ to
$\mathcal End_B(\mathcal G)$.
Compose it with (F1S33.5) and divide by the positive
integer $r$:
\[
 u_i:I/I^2\longrightarrow B,\qquad
          u_i(f)=\frac1r\operatorname{tr}_{\mathcal G}(a_{i,f}).
 \tag{F1S33.17}
\]
No scalarity assumption on $a_{i,f}$ is used.

The corresponding actual ideal is
\[
 I_{\varepsilon,i}
   =\{\,f+\varepsilon h:\ f\in I,\quad
                      \bar h=-u_i(f)\,\}
       \subset R[\varepsilon]/(\varepsilon^2).
 \tag{F1S33.18}
\]
$B$-linearity of $u_i$ proves that this is an ideal:
multiplication by $s+\varepsilon t$ sends the
coefficient condition to
$\overline{sh+tf}=-\bar s\,u_i(f)=-u_i(sf)$.
Its reduction is $I$, and
$I_{\varepsilon,i}\cap\varepsilon R=\varepsilon I$.
The quotient is flat over $A_\varepsilon$.
Indeed, if the class of $f+\varepsilon h$ is
annihilated by $\varepsilon$, then $f\in I$.
Modulo (F1S33.18), this class equals
$\varepsilon(h+\widetilde u_i(f))$, where
$\widetilde u_i(f)$ is any lift of $u_i(f)$ to $R$.
Thus the kernel of multiplication by $\varepsilon$
equals its image, and the central quotient is $B$.
This verifies flatness directly on the actual
quotient algebra.

\paragraph{4. The scalar twist disappears in the ideal gluing.}
Let $\gamma_{ij}$ denote the central reduction of
$\phi_{ij}$. The semilinearity in (F1S33.8) gives,
for every central $f\in I$,
\[
 \gamma_{ij}a_{i,f}\gamma_{ij}^{-1}
     =a_{j,f}+\kappa_{ij}(f)\,1_{\mathcal G}.
 \tag{F1S33.19}
\]
This follows by applying $\phi_{ij}$ to
$f\widetilde m$:
$\phi_{ij}(f\widetilde m)
 =(f+\varepsilon\kappa_{ij}(f))\phi_{ij}(\widetilde m)$,
and then using the actual identification in
(F1S33.14).
Conjugation invariance and trace of the identity
give
\[
       u_i(f)-u_j(f)=\kappa_{ij}(f)|_Y.
 \tag{F1S33.20}
\]
It follows that
$\rho_{ij}(I_{\varepsilon,i})=I_{\varepsilon,j}$:
the element $f-\varepsilon\widetilde u_i(f)$ is
sent to
$f+\varepsilon(\kappa_{ij}(f)-\widetilde u_i(f))$,
whose coefficient modulo $I$ is $-u_j(f)$.
The scalar in (F1S33.13) acts trivially under
conjugation. Thus it introduces no discrepancy
in (F1S33.19)--(F1S33.20).
The ideals glue in the honestly glued ambient
ring to an actual flat closed subscheme
$Y_\varepsilon\subset P_\kappa$.

The S24 normal obstruction $\eta_Y(\kappa)$ has
the following representative in the normal sheaf:
\[
  f\longmapsto\kappa_{ij}(f)|_Y,\qquad
  \mathcal Hom_Y(I/I^2,\mathcal O_Y).
\]
Equation (F1S33.20) displays it as a coboundary.
Using the actual S24 maps,
\[
 \eta_Y(\kappa)=0,\qquad
 \operatorname{ob}_E(\kappa)
       =d_2^{\,\mathrm{ideal}}j\,\eta_Y(\kappa)=0.
 \tag{F1S33.21}
\]
We have proved the required kernel inclusion
\[
 \ker\!\left(\operatorname{ev}_F\bigm|
                  _{D\oplus H^2(P,\mathcal O_P)}\right)
       \subseteq
 \ker\!\left((\kappa,\eta)\longmapsto
                               \operatorname{ob}_E(\kappa)\right).
 \tag{F1S33.22}
\]
This treats the scalar component $\eta$ throughout.

\paragraph{5. The characteristic-image quotient and comparison with S31.}
If two pairs $(\kappa,\eta)$ represent the same
element of $V_F$, their difference satisfies
(F1S33.7). Equation (F1S33.22) implies that their
images under $-r\,\operatorname{ob}_E$ agree.
This proves well-definedness and linearity of
(F1S33.3). It is onto because $r\ne0$ and
$W_E=\operatorname{ob}_E(D)$.
The quotient construction retains the actual
characteristic map of the given sheaf.

When $\mathcal G=i^*N$ for an actual ambient
vector bundle $N$ of rank $r$, the projection
formula gives the specified isomorphism
$i_*\mathcal G=E\otimes N$.
On the image $V_F$, S31 proves
\[
 p_{\mathrm{ns}}\operatorname{ptr}_N
                 \operatorname{ev}_F(\kappa+\eta)
               =-r\,\operatorname{ob}_E(\kappa).
 \tag{F1S33.23}
\]
Thus the map $\mathcal T_{\mathcal G}$ agrees
with that actual projected partial trace on
$V_F$. For general coherent $\mathcal G$, the
construction here supplies (F1S33.3) on this
specified image; it does not apply an
$\mathcal O_Y$-linear partial trace to an
arbitrary ambient self-Ext morphism.

\paragraph{6. The full character forces positive generic rank.}
The localized Chern-character map for the actual
closed immersion is
\[
\begin{split}
 K'_0(Y)&\longrightarrow K_0(D_{Y,\mathrm{perf}}(P))
       \xrightarrow{\ \operatorname{ch}(Y\to P,-)\cap[P]\ }
                     \mathrm{CH}_*(Y)_{\mathbf Q}\\
       &\xrightarrow{\ i_*\ }\mathrm{CH}_*(P)_{\mathbf Q}.
\end{split}
 \tag{F1S33.24}
\]
Its composite is the ordinary Chern character of
the exact pushforward, capped with $[P]$.
Stacks Proposition~42.57.1 (0FEW) proves that the
localized map respects support dimension; its
top component on an integral subscheme is the
fundamental cycle, with normalization supplied
by Lemma~42.55.4 (0FEH).
The regular ambient local ring in the
generic-point calculation is
$\mathcal O_{P,\xi}$.

For completeness the rank coefficient can be
seen on an actual exact sequence. Choose an
ambient ample line bundle $H$ and $n$ such that
$\mathcal G(nH)$ is generated by global sections.
Choose $r$ of these sections forming a generic
basis. They give an injection
$\mathcal O_Y(-nH)^{\oplus r}\to\mathcal G$:
its kernel is a torsion subsheaf of a locally
free sheaf on the integral $Y$, and is zero.
The cokernel $\mathcal C$ has support of
dimension at most three. For $r=0$, take
$\mathcal C=\mathcal G$.
Exact pushforward gives
\[
 0\longrightarrow E(-nH)^{\oplus r}
     \longrightarrow F\longrightarrow i_*\mathcal C
     \longrightarrow0.
 \tag{F1S33.25}
\]
The support-dimension property of (F1S33.24)
implies
$\operatorname{ch}_j(i_*\mathcal C)=0$ for $j\le4$.
S14 gives $\operatorname{ch}_j(E)=0$ for $j<4$
and $\operatorname{ch}_4(E)=[Y]$.
Apply the cycle-class morphism
$\operatorname{cl}:\mathrm{CH}^*(P)_{\mathbf Q}
 \to H^{2*}(P,\mathbf Q)$ to this calculation.
Additivity in (F1S33.25) therefore proves, in
rational cohomology,
\[
 \operatorname{ch}_j(F)=0\ (j<4),\qquad
 \operatorname{ch}_4(F)=r[Y]
        =r\left(w+\frac{\beta^4}{72}\right).
 \tag{F1S33.26}
\]
This equality is about the actual pushforward
(F1S33.2); no replacement of its sheaf by its
character is made in the obstruction proof.

In (F1S33.4), let $q_4$ be the coefficient of
$x^4$ in $Q$. Taking degree four and applying
the actual map
$\beta\wedge:H^8(P,\mathbf Q)\to H^{10}(P,\mathbf Q)$
gives
$r\beta^5/72=q_4\beta^5$,
since $\beta w=\beta t=0$.
The class $\beta^5$ is nonzero, as its product
with $\beta^3$ has positive integral.
It follows that $q_4=r/72$ and
\[
                         \lambda t=rw\ne0.
 \tag{F1S33.27}
\]
The generic rank is a nonnegative integer, so
(F1S33.27) proves $r>0$.

\paragraph{7. The full trace and every specified point lift.}
For the given $B$ define the actual HH component
\[
 a_B(\kappa)=\kappa+\kappa\mathbin{\lrcorner}B
       \in H^1(P,T_P)\oplus H^2(P,\mathcal O_P).
 \tag{F1S33.28}
\]
The polarized Weil tangent annihilates $\beta$
and the determinant classes. With
$D_\kappa=\kappa\mathbin{\lrcorner}(-)$,
the finite exponential identity gives
\[
 \bigl(D_\kappa+(D_\kappa B)\wedge\bigr)
       \bigl(e^{-B}(\lambda t+Q(\beta))\bigr)=0.
 \tag{F1S33.29}
\]
The entire BF characteristic--trace identity,
with the S31 sign convention, therefore gives
\[
 U_F:=\operatorname{ev}_F(a_B(D))
                         \subseteq\ker\sigma_F^2.
 \tag{F1S33.30}
\]
Restricting the already constructed map yields
the actual surjection
\[
 \left.-\frac1r\mathcal T_{\mathcal G}\right|_{U_F}:
     U_F\twoheadrightarrow W_E,\qquad
 \operatorname{ev}_F(a_B(\kappa))
                         \longmapsto\operatorname{ob}_E(\kappa).
 \tag{F1S33.31}
\]
In particular $\dim U_F\ge16-r_C$.

For every point functor $t_x^*(-)\otimes M$,
$M\in\operatorname{Pic}^0(P)$, its action on
$HH^2(P)$ is the identity: translations fix the
invariant-form description and $M$ has zero
Atiyah class. Naturality of the actual
characteristic action consequently fixes every
element of $U_F$ under any specified coherent
point lift on $F$.
Thus for its actual finite group $G$,
\[
 \dim\bigl(\ker\sigma_F^2
           \cap\operatorname{Ext}^2_P(F,F)^G\bigr)
                             \ge16-r_C\ge4.
 \tag{F1S33.32}
\]
The R05 identity-sector and induction/forgetting
morphisms identify this with the weak equivariant
trace kernel. The central object is the given
ordinary sheaf, so its negative self-Ext already
vanishes.

\paragraph{8. Concrete modifications and the retained scheme structure.}
For any actual zero-dimensional subscheme
$Z_0\subset Y$ of length $\ell$, form its ideal
sheaf $\mathcal I_{Z_0/Y}$ and the actual exact
sequence
\[
 0\longrightarrow i_*\mathcal I_{Z_0/Y}
    \longrightarrow E\longrightarrow i_*\mathcal O_{Z_0}
    \longrightarrow0.
 \tag{F1S33.33}
\]
Its generic rank on $Y$ is one. In the Chow
group, the localized top-cycle term is
$i_*[Z_0]$. Apply the actual cycle-class map
\[
 \operatorname{cl}:\mathrm{CH}^8(P)_{\mathbf Q}
                        \longrightarrow H^{16}(P,\mathbf Q).
\]
Every closed point has cohomology class of
integral one. Since $P$ is connected, its top
cohomology is one-dimensional, so
$\operatorname{cl}(i_*[Z_0])=\ell[p]$.
Thus the following equality is in rational
cohomology:
\[
 \operatorname{ch}(i_*\mathcal I_{Z_0/Y})
       =\operatorname{ch}(E)-\ell[p],
 \qquad
 [p]=\frac{\beta^8}{\int_P\beta^8}.
 \tag{F1S33.34}
\]
S14's full character remains of the required
form with $B=0$ and the same nonzero $w$.
Hence every such actual ideal modification has
the invariant kernel (F1S33.32).
When $Z_0$ is a point of the smooth locus, its
Koszul resolution on that locus and the
structure sheaf off the point also show that
$\mathcal I_{Z_0/Y}$ is perfect on $Y$.
The argument itself permits arbitrary coherent
$\mathcal G$, including torsion and singular
nonfree loci, under the exact quotient action
(F1S33.2).
The essential module identity used in
(F1S33.14)--(F1S33.16) is $I_YF=0$.
Sheaves annihilated only by a higher power of
$I_Y$, and derived objects with several nonzero
cohomology sheaves on $Y$, remain outside the
proved scope here.

\paragraph{Primary-source and framework scope.}
Toda's published Proposition~6.1 and its
Lemmas~6.2--6.4 were read in full through
printed pages 215--220, together with the
Section~4 ring gluing and Definition~4.1 of
twisted sheaves. The degree-one comparison
(F1S33.9) is the explicit formula on page 216.
The source's Proposition~5.3 is attributed there
to C\u ald\u araru without a supplied proof;
no fresh proof of that proposition is claimed,
and the needed degree-one comparison was
verified directly above.
The trace extension uses the complete displayed
proofs of Stacks 0AVB, 031S and 0AUW.
The support-dimension and top-cycle statements
use the complete proofs of Stacks 0FEW and
0FEH, applied with the smooth ambient $P$.
The remaining full characteristic--trace,
normal-obstruction and equivariant comparisons
are the exact BF/S24/S29/S31/R05 maps already
audited. In particular, the actual deformation
and ideal-gluing construction proves the
kernel inclusion before the linear quotient
(F1S33.3) is defined.
